\section{Fourier refinements at a small fixed index}
\label{ext:small-fourier}
At a small index, replacing a finite frequency integral by an integral
over the whole half-line can lose much of the available margin. The
following estimates retain the finite interval and the actual integer
power. They use the moments of the same supported marginal throughout.

\subsection{A finite-interval compensated estimate}
Let $n>2$ be an integer, $\E X=0$, $\E X^2=1$, $|X|\le R\le6$,
$\gamma=\E X^3$, $\mu_j=\E X^j$, and $f(t)=\E e^{itX}$.
Assume $\mu_4\le M\le111/10$ and $|\gamma|\le g\le37/20$.
The coefficient and moment bounds proved in the formula atlas give
\[
 |a|\le A_*,\quad |z|\le Z,\quad \mu_6\le M_6,
 \qquad a=(3-\mu_4)/24,\quad z=(10\gamma-\mu_5)/120.
\]
Define
\[
 E=\frac{M_6}{720}+\frac{g^2}{72}
       +\frac{(1/2+g^2/144)^3}{5},\qquad
 B=A_*+Z/2+E/4,
\]
\[
 a_n=(n-1)/2,\qquad b_n=(n-2)(45-M)/96,\qquad
 I_j(A)=\int_0^{1/2}t^j e^{-At^2}\,dt.
\]
All these quantities are nonnegative, and the two damping rates are
positive. Put
\begin{align*}
 L_0={}&\sqrt n\left\{nZI_2(a_n)
       +\frac{n(n-1)A_*g}{6}I_4(a_n)
       +\frac{n^3g^3}{1296}I_6(n/2)\right\},\\
 L_H={}&\sqrt n\left\{nA_*I_2(a_n)
                              +\frac{n^2g^2}{72}I_4(n/2)\right\},\\
 L_2={}&n\left\{nEI_3(a_n)+\frac{n(n-1)B^2}{2}I_5(b_n)\right\}.
\end{align*}
For every real raw threshold $H$, the complete low-frequency estimate is
\begin{equation}\label{ext:finite-low}
 \int_0^{1/2}\frac{\left|\Im\left[e^{-itH}
 \{f(t)^n-e^{-nt^2/2}(1-in\gamma t^3/6)\}\right]\right|}{t^3}\,dt
 \le\frac{L_0+L_H|H|}{\sqrt n}+\frac{L_2}{n}.
\end{equation}

To prove it, set $h(t)=e^{-t^2/2-i\gamma t^3/6}$. Taylor's formula gives
\[
 f-h=-at^4-izt^5+r_6,\qquad |r_6|\le Et^6,
 \qquad |f-h|\le Bt^4 \quad(0\le t\le1/2).
\]
Indeed, for $w=-t^2/2-i\gamma t^3/6$, one has
$|w|\le t^2(1/2+g^2/144)<1/6$. The quadratic exponential
polynomial contributes at most $g^2t^6/72$ beyond degree five; its
remaining error is at most $|w|^3e^{|w|}/6\le |w|^3/5$.
The characteristic-function remainder is at most $M_6t^6/720$.
Symmetrization and the global fourth-degree cosine upper bound give
\[
 |f(t)|^2\le1-t^2+(M+3)t^4/12
           \le1-(45-M)t^2/48.
\]
Thus $|f|$, $|h|$, and every point of their complex line segment are
at most $e^{-(45-M)t^2/96}$. The integral Taylor remainder for the
polynomial $z\mapsto z^n$ consequently yields
\[
 f^n-h^n=n(f-h)h^{n-1}+R_2,\qquad
 |R_2|\le\tfrac12 n(n-1)B^2t^8e^{-b_nt^2}.
\]
The $r_6$ and $R_2$ integrals give $L_2/n$. For the leading terms use
\[
 |\Im[e^{-itH}(-at^4-izt^5)h^{n-1}]|
 \le e^{-a_nt^2}\{(Z+A_*|H|)t^5+(n-1)A_*gt^7/6\}.
\]
Finally, for $u=n\gamma t^3/6$, the rotated imaginary part of
$e^{-nt^2/2}(e^{-iu}-1+iu)$ is bounded by
$e^{-nt^2/2}(|u|^3/6+|H|tu^2/2)$. Integration proves
\eqref{ext:finite-low}, including its continuous endpoint at zero.

For the small indices the quantities $I_j(A)$ are enclosed directly.
The degree-20 and degree-21 Taylor polynomials of $e^{-x}$ are,
respectively, upper and lower bounds on all $x\ge0$, by the sign of
the integral remainder. Their integrated values are the rational sums
\[
 \sum_{k=0}^{m}\frac{(-A)^k}{k!(j+2k+1)2^{j+2k+1}},
 \qquad m=20,21.
\]
The selected routines require $0\le A\le8$ and $0\le j\le16$;
the rates above meet those gates for $5\le n\le9$.
Every coefficient multiplying an integral is nonnegative. When two
complete triples $(L_0,L_H,L_2)$ are available, either may be selected
by a rule independent of $H$. Selecting different coefficients from
different triples is not justified by this argument.

\subsection{Containing rectangles and exact finite powers}
The supported fourth-moment bound also improves the real coordinate.
If $\rho\in[\rho_l,\rho_u]$, $|X|\le R$, and $\mu_4\le M_4$, then
\[
 \Re f(t)\le1-t^2/2+M_4t^4/24.
\]
If $R^2t^2\le30$, a second lower endpoint is
\[
 \Re f(t)\ge1-t^2/2
              +\rho_l^2t^4(1-R^2t^2/30)/24.
\]
These follow from the global fourth-degree upper and sixth-degree
lower Taylor polynomials of cosine, $\mu_6\le R^2\mu_4$, and
$\mu_4\ge\rho^2$ by Cauchy's inequality. The nonnegative coefficient
gate is needed before substituting $\rho_l^2$. Intersect these real
intervals with the anchored bounds of Lemma~\ref{ext:fourier-anchor}.

\begin{lemma}[Two corners of a complex-power rectangle]
\label{ext:two-corners}
Let $n\ge2$ be an integer, $a\ge0$, $0<x_0\le x_1$ and $Y\ge0$.
If
\[
 (n-2)Y\le3x_0/2,\qquad (2n-1)x_0^n\ge(n-1)a,
\]
then for $x\in[x_0,x_1]$ and $|y|\le Y$,
\[
 |(x+iy)^n-a|^2
 \le\max\{|(x_0+iY)^n-a|^2,|(x_1+iY)^n-a|^2\}.
\]
\end{lemma}
\begin{proof}
Write $z=x+iy$, $r=|z|$ and $F(x,y)=|z^n-a|^2$.
Since $|\arg z|\le Y/x_0$, the first gate and $3/2<\pi/2$ imply
$\Re z^{n-2}\ge0$. Direct differentiation gives
\begin{align*}
 F_{yy}&=2nr^{2n-4}\{r^2+2(n-1)y^2\}
                         +2an(n-1)\Re z^{n-2}\ge0,\\
 F_{xx}&=2nr^{2n-4}\{r^2+2(n-1)x^2\}
                         -2an(n-1)\Re z^{n-2}.
\end{align*}
Divide the latter by $2nr^{n-2}>0$. Its first term is
$r^{n-2}((2n-1)x^2+y^2)\ge(2n-1)x_0^n$; the subtracted
term is at most $a(n-1)$. Thus $F$ is convex in each coordinate.
Apply convexity first in $y$ and then in $x$, and use
$F(x,Y)=F(x,-Y)$.
\end{proof}

For Fourier smoothing take $a=e^{-nt^2/2}$ and any proved containing
rectangle for $f(t)$. Both gates must hold throughout the frequency
cell. If either fails, the geometric-sum estimate remains available.
The maximum of the two pointwise corner bounds is enclosed before
integration; the maximum of their separate integrals would not suffice.
Complex powers are finite polynomial recurrences, so zero or equal
bases create no division. A fourth-Taylor endpoint is kept fixed on
each integration cell when differentiating its smooth majorant.

Cancellation can be retained when computing the real endpoints.
For example, on $0\le u\le1$,
\[
 1-u^2/2+M_4u^4/24-e^{-u^2/2}
 =\frac{M_4-3}{24}u^4
   +\sum_{k=3}^{10}\frac{(-1)^{k+1}u^{2k}}{2^kk!}+R(u).
\]
For derivative order $j=0,1,2$ with $u=ct$, the absolute remainder
is at most $2(22)_j u_{\max}^{22-j}c^j/(2^{11}11!)$; consecutive
absolute tail terms have ratio below $1/2$ on this domain.
Alternatively, at higher frequencies enclose each explicit smooth
residual $d(u)=\cos u-e^{-u^2/2}$ or
$d(u)+(\rho_u-1)u^3/6$ by its midpoint value plus the half-width
times an enclosure of its derivative on the entire cell.
Intersecting two such valid enclosures preserves their common function.
These operations differentiate scalar envelopes, not the law's CDF.

\section{A five-summand corner without a support assumption}
\label{ext:gapfree5}
The last small moment rectangle can be treated directly by conditioning
on signs. This avoids contact and Fourier estimates in that rectangle.
Let $\beta=\E[X|X|]$, $p=\Pp(X>0)$, $q=1-p$,
$m=\E X_+=\E X_-$ and $a_\pm=(1\pm\beta)/2$.
Zero atoms belong to the nonpositive class. Signwise Cauchy gives
\[
 m\ge\frac{1+\beta^2}{2\rho},\qquad
 \frac{m^2}{a_+}\le p\le1-\frac{m^2}{a_-},\qquad m^2\le pq.
\]
Set $h=pq/m$, $Y=p+hX$, $B=\ind_{\{X>0\}}$, and $U=Y-B$.
Both conditional means of $U$ vanish, with variances
\[
 v_0=p^2(qa_-/m^2-1),\qquad v_1=q^2(pa_+/m^2-1).
\]
Given $\sum_{i=1}^5B_i=j$, the noises are independent with variance
$V_j=(5-j)v_0+jv_1$. Write $b_j=\binom5j p^jq^{5-j}$ and define
\[
 T(V,d)=\begin{cases}1,&d\le0,\\V/(V+d^2),&d>0.\end{cases}
\]
Conditional Cantelli and the exact normal coordinate yield
\begin{equation}\label{ext:gapfree-bound}
 \frac{\sqrt5}{\rho}\{\Pp(S_5\le t)-\Phi(t/\sqrt5)\}
 \le\frac{\sqrt5}{\rho}
 \left\{\sum_{j=0}^5 b_jT(V_j,j-z)
          -\Phi\left(\frac{(z-5p)m}{\sqrt5pq}\right)\right\},
 \qquad z=5p+ht.
\end{equation}
The value $T(0,0)=1$ is essential: a deterministic sum at the threshold
belongs to the closed event.

\begin{proposition}[The positive corner]
\label{ext:gapfree5-result}
For every standardized iid marginal satisfying
\[
 \frac{477}{400}\le\rho\le\frac65,\qquad
 \frac{17}{50}\le\beta\le\frac{35139006743001}{10^{14}},
\]
the positive closed-CDF score in \eqref{ext:gapfree-bound} is strictly
below $4097/10000<\CE$ at every real threshold.
\end{proposition}
\begin{proof}
Here are the enclosures defining the finite proof. On a closed
$\beta$ slab $[b_l,b_u]$ put $m_l=(1+b_l^2)/(2(6/5))$.
The moment inequalities enclose $p$ in
$[2m_l^2/(1+b_u),\,1-2m_l^2/(1-b_l)]$.
On a closed $p,z$ box $[l,u]\times[a,b]\subset(0,1/2)\times[0,5]$
use
\begin{align*}
 v_0&\le\max\{0,u^2((1-l)(1-b_l)/(2m_l^2)-1)\},\\
 v_1&\le\max\{0,(1-l)^2(u(1+b_u)/(2m_l^2)-1)\}.
\end{align*}
Each binomial mass is maximized at the projection of $j/5$ onto
$[l,u]$, and each Cantelli factor is maximized at $z=b$ and the
variance upper bound. Clip the sum of the six upper bounds at one.
If $a-5u\ge0$, the normal argument is at least
$(a-5u)m_l/(\sqrt5\,u(1-u))$.
Otherwise it is at least $(a-5u)/\sqrt{5u(1-u)}$: use
$m\le\sqrt{pq}$ on the negative part and the fact that
$(z-5p)/\sqrt{pq}$ decreases in $p$ for $z\ge0$, $p<1/2$.
Its derivative has numerator $-[5p+z(1-2p)]/2$.
Outward square roots and a rational Gaussian lower enclosure therefore
bound the complete box. Take the nonnegative part of the CDF difference
before multiplying by $\sqrt5/\rho_l$.

The twelve closed $\beta$ slabs have endpoints
\[
 17/50,\quad 341/1000,\ldots,351/1000,\quad
 35139006743001/10^{14}.
\]
Together with the five intervals $z\in[k,k+1]$, $0\le k<5$, these
give sixty roots. The supplied binary subdivisions have 1,027 leaves.
Exact evaluation on every closed leaf gives a maximum upper bound
\[
 204848338979/500000000000<4097/10000.
\]
The independent evaluator uses a 100-term integrated Gaussian Taylor
polynomial with its explicit remainder, rational Machin bounds and
integer square-root enclosures. Complete root order, both children at
every split, and all leaf bounds are checked; twenty-one negative
controls reject missing cells, false bounds and malformed data.

Both omitted rays have direct bounds. Use the full rectangle's $m_l$
and enclosing $[p_l,p_u]$, put $v_u=p_u(1-p_u)$ and
$s_u=v_u^2/m_l^2-v_u$. On the feasible domain $m^2\le v$,
$v^2/m^2-v$ increases with $v$ and decreases with $m$. Thus
$V_j\le\max\{j/p_l,(5-j)/(1-p_u)\}s_u$.
For $z\le0$, discard the nonnegative Gaussian CDF and evaluate the
Cantelli upper sum at zero, using binomial upper masses
$\binom5j p_u^j(1-p_l)^{5-j}$. The resulting score is below $331/1000$.
For $z\ge5$, the CDF is at most one and normal Cantelli gives the
score upper bound
\[
 \frac{\sqrt5}{\rho_l}\frac1{1+5m_l^2/p_u^2}<187/1000.
\]
These exact rational comparisons and the sixty roots cover all real $z$.
The proposition concerns the stated positive branch; it is not silently
extended to the opposite $\beta$ rectangle by reflection.
\end{proof}
